% from ega3-1-fr.tex
\oldpage[III]{134}
$X'=X\times_YY'$, $f'=f\times1_{Y'}$~; désignons par $\eta$, $\eta'$,
$\xi$ les points génériques de $Y$, $Y'$ et $X$ respectivement, de sorte
que $f(\xi)=\eta$ et $h^{-1}(\eta)=\{\eta'\}$~; en outre,
$\boldsymbol{k}(\eta)=\boldsymbol{k}(\eta')=\mathrm R(Y)$, donc $f'^{-1}(\eta')$ est isomorphe à
$f^{-1}(\eta)$ (\textbf{I}, 3.6.4) et en particulier, puisque $\xi$ est le point
générique de $f^{-1}(\eta)$ (\textbf{0}\textsubscript{I}, 2.1.8),
$f'^{-1}(\eta')$ a un seul point générique. Mais par hypothèse, toute
composante irréductible de $X'$ a son point générique dans
$f'^{-1}(\eta')$, donc $X'$ est nécessairement irréductible, son point
générique $\xi'$ est le point générique de $f'^{-1}(\eta')$ et on a
$\boldsymbol{k}(\xi')=\boldsymbol{k}(\xi)$.
Posons $X''=X'_{\mathrm{red}}$, $f''=f'_{\mathrm{red}}$~; $X''$ est donc
intègre et noethérien, $f''$ est propre (\textbf{II}, 5.4.6) et les espaces
sous-jacents aux fibres $f'^{-1}(y')$ et $f''^{-1}(y')$ sont les mêmes~; en
outre, $\mathrm R(X'')=\boldsymbol{k}(\xi')=\mathrm R(X)$, donc $f''$ vérifie les
hypothèses de (\hyperref[III.4.3.8-fr]{4.3.8}), et $f''^{-1}(y')$ est
géométriquement connexe. Soit maintenant $k$ une extension quelconque de
$\boldsymbol{k}(y)$~; il existe une extension $k_1$ de $\boldsymbol{k}(y)$ dont $\boldsymbol{k}(y')$ et $k$ peuvent
être considérés comme des sous-extensions (Bourbaki, \emph{Alg.}, chap. V,
\S~4, prop. 2). Par hypothèse,
$f''^{-1}(y')\times_{Y'}\operatorname{Spec}(k_1)$ est connexe, et il a même
préschéma réduit que
$f'^{-1}(y')\times_{Y'}\operatorname{Spec}(k_1)$ (\textbf{I}, 5.1.8), donc ce dernier
est connexe, et comme il est isomorphe à
$f^{-1}(y)\times_Y\operatorname{Spec}(k_1)$ (\textbf{I}, 3.6.4), on en conclut que ce
dernier est connexe~; \emph{a fortiori}, il en est de même de
$f^{-1}(y)\times_Y\operatorname{Spec}(k)$ d'après la remarque du début de
(\hyperref[III.4.3.4-fr]{4.3.4}), ce qui achève la démonstration.
\end{proof}

\begin{env}[4.3.11]
\label{III.4.3.11-fr}
\emph{Remarques} (4.3.11). ---
\begin{enumerate}
\item[(i)] Le raisonnement précédent est dû en substance à Zariski [20], à
cela près qu'il peut prendre pour $A$ la clôture intégrale de
$\mathcal O_y$, celle-ci étant un anneau noethérien pour les anneaux locaux
de la géométrie algébrique classique. D'autre part, Zariski prouve que si
$Y$ est la variété de Chow d'un espace projectif $\mathbf P_k^r$ sur un
corps $k$, et si $X$ est la partie fermée de $\mathbf P_k^r\times_kY$ qui
définit la correspondance de Chow entre $\mathbf P_k^r$ et $Y$, alors la
projection $X\to Y$ est un morphisme universellement ouvert
(\emph{loc. cit.}, lemme de la p. 82). Il semble bien que ce soit la seule
propriété formelle des «~coordonnées de Chow~» dont on se soit servi dans
certaines applications~; il y a par suite intérêt dans une telle situation,
à substituer le langage~: fibres d'un morphisme propre (éventuellement
supposé universellement ouvert ou soumis à d'autres restrictions analogues
de régularité locale) au langage~: spécialisation de cycles dans l'espace
projectif.

\item[(ii)] Au chap. IV, nous verrons qu'un morphisme universellement ouvert
$f:X\to Y$ peut encore être défini de la façon suivante (qui justifie la
terminologie)~: pour \emph{tout} morphisme $Z\to Y$, le morphisme
$f_{(Z)}:X_{(Z)}\to Z$ est \emph{ouvert}. On peut montrer en outre que si $f$
vérifie les hypothèses de (\hyperref[III.4.3.10-fr]{4.3.10}), alors, si
$y,y'$ sont deux points de $Y$ tels que $y$ soit une \emph{spécialisation} de $y'$,
le nombre géométrique de composantes connexes de $f^{-1}(y)$ est au plus
égal à celui des composantes connexes de $f^{-1}(y')$.
\end{enumerate}
\end{env}

\begin{corollary}[4.3.12]
\label{III.4.3.12-fr}
Sous les hypothèses de \textup{(\hyperref[III.4.3.5-fr]{4.3.5})}, supposons en outre
$\mathrm R(Y)$ algébriquement fermé dans $\mathrm R(X)$, et soit $y$ un
point normal de $Y$. Alors $f^{-1}(y)$ est géométriquement connexe, et il
existe un voisinage ouvert $U$ de $y$ dans $Y$ tel que
$f_*(\mathcal O_X|f^{-1}(U))$ soit isomorphe à $\mathcal O_Y|U$. Si plus
particulièrement, on suppose $Y$ normal (et $\mathrm R(Y)$ algébriquement
fermé dans $\mathrm R(X)$) $f_*(\mathcal O_X)$ est isomorphe à
$\mathcal O_Y$.
\end{corollary}

\begin{proof}
La première assertion relative à $f^{-1}(y)$ est un cas particulier de
(\hyperref[III.4.3.8-fr]{4.3.8}). On en
